\section{Kiến trúc tăng tốc CNN gắn với RISC-V}
\subsection{RoCC, DMA và tổ chức GEMM}
Nghiên cứu của He là đối chiếu chính về đường tích hợp CPU--accelerator. Phần kiến trúc mô tả Rocket Core kết nối bộ tăng tốc qua RoCC; bộ tăng tốc dùng DMA để đọc/ghi feature map qua L2 cache và DRAM. Khối điều khiển, khối tính toán và khối lưu trữ cục bộ cho phép phân biệt việc phát lệnh của CPU với di chuyển dữ liệu của phần cứng~\cite[Section~3, Figure~4]{he2021design}. Khác với cơ chế này, trong MyTinyMobileNetSoC, CPU phải đọc tensor rồi ghi toán hạng qua MMIO; ngoại vi không có đường DMA tới RAM.

Khối tính toán của He kết hợp GEMM với mảng systolic giữ trọng số tại các phần tử xử lý (weight-stationary). Các PE hợp thành tile; thanh ghi giữa các tile tổ chức đường ống. Trọng số được nạp trước rồi dùng với luồng đầu vào, trong khi chia khối phép nhân ma trận nhằm giảm truy cập lặp~\cite[Sections~2.2--3, Figure~5]{he2021design}. Phạm vi mô tả tập trung vào tăng tốc tích chập bằng tính toán ma trận; sơ đồ có khối activation/pooling, nhưng bài không mô tả đủ để xác nhận một đường depthwise chuyên biệt hoặc toàn bộ hậu xử lý giống firmware của luận văn.

Đặt cạnh thiết kế này, Conv/Depthwise/Pointwise của MyTinyMobileNetSoC biểu diễn các phép tính theo cấu trúc toán tử thay vì một mảng GEMM chung. Firmware quyết định khi nào giữ patch, nạp trọng số và sử dụng kết quả cho từng shape. Phương án ma trận có thể dùng chung lõi tính toán cho nhiều lớp, đổi lại yêu cầu ánh xạ tensor và cấp dữ liệu đủ đều. Phương án chuyên biệt giảm nhu cầu ánh xạ qua ma trận trong đường chạy hiện tại, nhưng có phạm vi shape và các nhánh fallback phải kiểm chứng. Hai cách tổ chức đặt sự phức tạp ở những vị trí khác nhau.

\subsection{Giới hạn bằng chứng và ý nghĩa đối chiếu}
Phần thực nghiệm của He sử dụng YOLOv3-tiny và báo cáo thời gian xử lý trên nền tảng FPGA~\cite[Section~4]{he2021design}. Tuy nhiên, sơ đồ kiến trúc ghi Rocket Core còn phần kết luận nêu BOOM hai lõi. Vì thế, chương chỉ dùng Rocket/RoCC để chỉ phương án kết nối được mô tả, không coi đây là cấu hình host đã được xác nhận nhất quán cho toàn bộ số đo. Bài cũng không công bố rõ độ rộng toán hạng và tích lũy của cấu hình đánh giá, nên không gán INT8/INT32 cho thiết kế này.

Sự khác nhau về workload, host và hệ thống nhớ không cho phép xếp hạng từ FPS được báo cáo. Giá trị đối chiếu nằm ở mức tự chủ: với DMA, CPU có thể giao tác vụ có kích thước lớn hơn cho accelerator; với MMIO do CPU cấp dữ liệu, chi phí mỗi lần gọi và lượng dữ liệu hữu ích mỗi giao dịch trở thành mục tiêu tối ưu trực tiếp. Đây là cơ sở để Chương~\ref{chap:kientruc} làm rõ ranh giới giao diện và Chương~\ref{chap:toiuu} đánh giá cả phần truyền dữ liệu của MyTinyMobileNetSoC.
